\documentclass[10pt,a4paper]{article}
\usepackage{xeCJK}
\usepackage{fontspec}
\usepackage{xcolor}
\usepackage{fancyhdr}
\usepackage{booktabs}
\usepackage{array}
\usepackage{amssymb}
\usepackage{geometry}
\usepackage[protrusion=true]{microtype}
\geometry{margin=1.7cm,top=1.9cm,headheight=15pt}
\linespread{1.0}

\setmainfont{Baskerville}
\setCJKmainfont{Songti SC}
\newfontfamily{\starfont}{Songti SC}

\definecolor{wrongred}{RGB}{198,40,40}
\definecolor{wrongbg}{RGB}{253,235,233}
\definecolor{rulegray}{RGB}{200,200,200}

\setlength{\emergencystretch}{3em}
\setlength{\fboxsep}{3pt}
\setlength{\parindent}{0pt}
\newcommand{\STAR}{\textcolor{wrongred}{\starfont ★}}
\newcommand{\W}[1]{\textcolor{wrongred}{#1}}
\newcommand{\A}[1]{\textbf{#1}}
\newcommand{\OK}{\textcolor{wrongred}{\ensuremath{\checkmark}}}
\newcommand{\NO}{\textcolor{wrongred}{$\times$}}
\newcommand{\UNC}{\textcolor{wrongred}{?}}
\newcommand{\checkmarkbox}{\fbox{\rule{0pt}{0.8ex}\rule{0.8ex}{0pt}}}
\newcolumntype{L}[1]{>{\raggedright\arraybackslash}p{#1}}

\pagestyle{fancy}
\fancyhf{}
\fancyhead[L]{\footnotesize 英语短语默写 小蓝 P41--50（错题订正）}
\fancyhead[R]{\footnotesize 赵文瀚}
\fancyfoot[C]{\thepage}
\renewcommand{\headrulewidth}{0.4pt}
\renewcommand{\footrulewidth}{0pt}

\begin{document}

\begin{center}
{\LARGE\bfseries 英语短语默写 小蓝 P41--50}\\[0.25em]
{\large 错题订正与复盘}\\[0.35em]
{\footnotesize 2029 届高一英语 · 原卷 25 条 · 2026-10}
\end{center}

\vspace{-0.2em}
\noindent{\color{rulegray}\rule{\textwidth}{0.8pt}}

\vspace{0.5em}
\noindent\fcolorbox{black!15}{black!4}{%
  \parbox{\dimexpr\linewidth-2\fboxsep-2\fboxrule\relax}{%
    \small 来源：小蓝短语默写 P41--50（共 25 条）\qquad 姓名：赵文瀚\qquad 重做日期：\underline{\hspace{2.4cm}}}}

\vspace{0.55em}
\noindent{\footnotesize 说明：按原卷题号列全 25 条。\W{红色}＝可辨的学生原答，\A{加粗}＝参考表达（并非原卷答案）；\OK 表示初核可接受，\NO 表示初核不符，\UNC 表示原答或红笔判定无法确证。原卷右下圈记“10”，含义未注明，不据此推算对错数。}

\vspace{0.35em}
\renewcommand{\arraystretch}{0.90}
\noindent{\scriptsize\begin{tabular}{@{}L{0.55cm}@{\hspace{0.05cm}}L{1.45cm}@{\hspace{0.05cm}}L{1.25cm}@{\hspace{0.05cm}}L{4.0cm}@{\hspace{0.05cm}}L{4.7cm}@{\hspace{0.05cm}}L{0.62cm}@{\hspace{0.05cm}}L{1.15cm}@{\hspace{0.05cm}}L{0.55cm}@{}}
\toprule
\textbf{题号} & \textbf{中文短语} & \textbf{提示} & \textbf{学生原答} & \textbf{参考表达} & \textbf{核对} & \textbf{初核问题} & \textbf{二刷} \\
\midrule
1 & 回顾 \STAR & — & \W{get back to}（旁见改写 look back…；全答不清） & \A{look back (on sth.)} & \UNC & 涂改待核 & \checkmarkbox \\
2 & 羽毛球场 & — & badminton court & \A{badminton court} & \OK & — & \checkmarkbox \\
3 & 落伍 \STAR & — & \W{get out of fashion}（有涂改） & \A{go / be out of fashion} & \NO & 搭配 & \checkmarkbox \\
4 & 打钟，按铃 & — & ring the bell & \A{ring the bell} & \OK & — & \checkmarkbox \\
5 & 廉价销售 \STAR & — & \W{sell … at a bargain sale}（多处涂改） & \A{sell at a bargain price} & \UNC & 涂改待核 & \checkmarkbox \\
6 & 向某人乞求某物 & — & beg sb. for sth. & \A{beg sb. for sth.} & \OK & — & \checkmarkbox \\
7 & 铺床 & — & make the bed & \A{make the bed} & \OK & — & \checkmarkbox \\
8 & 境况好转 & — & get better（旁见 be better off） & \A{get better / be better off} & \UNC & 旁注待核 & \checkmarkbox \\
9 & 规矩点！ & — & behave yourself! & \A{behave yourself!} & \OK & — & \checkmarkbox \\
10 & 难以描述 \STAR & — & \W{hard to describe beyond words} & \A{hard to describe} & \NO & 重复表达 & \checkmarkbox \\
11 & 束紧裤带 \STAR & — & \W{tighten the belt} & \A{tighten one's belt} & \NO & 固定表达 & \checkmarkbox \\
12 & 畅销书 & — & best-seller & \A{best-seller} & \OK & — & \checkmarkbox \\
13 & 破坏微妙的自然平衡 \STAR & — & \W{break the balance of nature}（疑有补词） & \A{break the delicate balance of nature} & \NO & 漏译修饰语 & \checkmarkbox \\
14 & 禁止某人做某事 & ban & ban sb. from doing sth. & \A{ban sb. from doing sth.} & \OK & — & \checkmarkbox \\
15 & 她刚到就得走了。 \STAR & Barely & \W{Barely she got …（后续涂改）} & \A{Barely had she arrived when she had to leave.} & \UNC & 语序／涂改待核 & \checkmarkbox \\
16 & 良好的开端是成功的一半。 & — & Well begun is half done. & \A{Well begun is half done.} & \OK & — & \checkmarkbox \\
17 & 不怕迟到只怕不做。 \STAR & — & \W{Don't be afraid of being late, but don't do it.}；红笔改写 Better late than never. & \A{Better late than never.} & \NO & 谚语误用 & \checkmarkbox \\
18 & 这本书我看不懂。 \STAR & beyond & This book is beyond my knowledge. & \A{This book is beyond me / beyond my understanding.} & \UNC & 搭配待核 & \checkmarkbox \\
19 & 不能忍受她总是迟到。 & bear & can't bear her being late & \A{can't bear her being late all the time} & \UNC & 漏译频度待核 & \checkmarkbox \\
20 & 他还要很长时间才能回来。 & It, before & It'll be long before he comes back. & \A{It'll be a long time before he comes back.} & \UNC & 涂改待核 & \checkmarkbox \\
21 & 新产品是基于传统草药研制的。 & base & The new product is based on traditional herbal medicine. & \A{The new product is based on traditional herbal medicine.} & \OK & — & \checkmarkbox \\
22 & 他按过门铃后好一会儿才开门。 \STAR & It, before & \W{It was long before he knocked the door.（涂改）} & \A{It was a long time before he opened the door.} & \NO & 语义／搭配 & \checkmarkbox \\
23 & 请允许我代表在座各位说你退休后生活愉快。 \STAR & behalf, present & \W{Please allow me to bless you a nice retired life on behalf of everyone present.（涂改）} & \A{Please allow me, on behalf of everyone present, to wish you a happy retirement.} & \NO & 词义／搭配 & \checkmarkbox \\
24 & 多读书对精神成长有益。 \STAR & \shortstack[l]{beneficial,\\expose} & \W{It's beneficial for you to read books often.} & \A{It's beneficial to expose yourself to more reading.} & \NO & 未用提示词／漏译 & \checkmarkbox \\
25 & 把我的字典递给我好吗？ & belong & Can you bring (me?) the dictionary which belongs to me? & \A{Can you pass me the dictionary that belongs to me?} & \UNC & 涂改待核 & \checkmarkbox \\
\bottomrule
\end{tabular}}

\vspace{0.2em}
\noindent{\scriptsize\textbf{原卷旁注}：日期未印，按收件归档月份记为 2026-10。第 1、5、15、25 题原答有涂改，完整作答无法可靠辨认；第 8、20 题原答可辨，红笔判定含义不明；第 18、19 题原答可辨，待核的是表达与中文原意是否相符（第 19 题是否体现“总是”）。以上标 \UNC 待复核。第 17 题旁可辨红笔改写“Better late than never.”。表中参考表达为整理时拟列，并非原卷所附答案；右下圈记“10”未擅自解释为分数或错题数。原图见本套归档扫描件。}

\newpage
\begin{center}
{\Large\bfseries 错因归类与复盘}
\end{center}
\vspace{-0.4em}
\noindent{\color{rulegray}\rule{\textwidth}{0.8pt}}

\vspace{0.5em}
\noindent 说明：按本次可明确核出的表达问题暂列；待核题目不计入条数，红笔批改痕迹需对照原图复核。

\vspace{0.5em}
\renewcommand{\arraystretch}{1.25}
\begin{center}
\begin{tabular}{@{}p{4.6cm} p{10.3cm} c@{}}
\toprule
\textbf{错误类型} & \textbf{题号} & \textbf{条数} \\
\midrule
固定搭配／用词 & {3, 11, 22, 23} & 4 \\
重复表达／漏译／未用提示词 & {10, 13, 24} & 3 \\
谚语／句式 & {17} & 1 \\
\midrule
\textbf{明确不符小计} & \textbf{3, 10, 11, 13, 17, 22, 23, 24} & \textbf{8} \\
\textbf{原答或红笔待核} & \textbf{1, 5, 8, 15, 18, 19, 20, 25} & \textbf{8} \\
\bottomrule
\end{tabular}
\end{center}

\vspace{0.5em}
\noindent\textbf{关联知识点标签}
\vspace{0.25em}

\noindent\begin{tabular}{@{}p{8.2cm}@{\hspace{0.4cm}}p{8.2cm}@{}}
\begin{minipage}[t]{\linewidth}
\small
\STAR\ \textbf{go / be out of fashion}：表示“落伍、过时”\\[0.2em]
\STAR\ \textbf{tighten one's belt}：固定表达，表示“束紧裤带”\\[0.2em]
\STAR\ \textbf{delicate balance of nature}：微妙的自然平衡\\[0.2em]
\STAR\ \textbf{Barely ... when ...}：句首 Barely 时主句倒装\\[0.2em]
\STAR\ \textbf{Better late than never}：迟做总比不做好
\end{minipage} &
\begin{minipage}[t]{\linewidth}
\small
\STAR\ \textbf{beyond me / beyond my understanding}：超出理解范围\\[0.2em]
\STAR\ \textbf{bear}：can't bear sb. doing sth.\\[0.2em]
\STAR\ \textbf{It is beneficial to expose oneself to ...}：有益于接触／阅读……\\[0.2em]
\STAR\ \textbf{on behalf of ...; wish sb. a happy retirement}：代表……祝某人退休愉快
\end{minipage} \\
\end{tabular}

\vspace{0.65em}
\noindent\fcolorbox{black!15}{black!4}{%
  \parbox{\dimexpr\linewidth-2\fboxsep-2\fboxrule\relax}{%
    \textbf{复盘记录}\\[0.35em]
    最薄弱的板块：\underline{\hspace{9.0cm}}\\[0.55em]
    主要错误原因：\checkmarkbox 固定搭配不牢\quad \checkmarkbox 介词记混\quad \checkmarkbox 漏词/多词\quad \checkmarkbox 用词/拼写\quad \checkmarkbox 空白未答\\[0.55em]
    本次复习的改进措施：\underline{\hspace{10.1cm}}\\[0.55em]
    \hspace{3.2em}\underline{\hspace{10.1cm}}\\[0.55em]
    下次复习日期：\underline{\hspace{2.2cm}}\qquad
    当前掌握度：\checkmarkbox 仍薄弱\quad \checkmarkbox 基本掌握\quad \checkmarkbox 已掌握
  }}

\end{document}
